\documentclass[10pt,a4paper]{amsart}
\usepackage[margin=19mm]{geometry}
\usepackage{amsmath,amssymb,mathtools}
\usepackage[scr=rsfs,cal=euler]{mathalfa}
\usepackage{xfrac}
\usepackage[unicode,hidelinks,bookmarks=false]{hyperref}
\newcommand{\ie}{i.e.\ }
\newcommand{\cf}{cf.\ }
\newcommand{\resp}{respectively\ }
\newcommand{\Subsection}{\subsection}
\providecommand{\nobreakdash}{\nobreak\hbox{-}}
\setlength{\emergencystretch}{2em}
\multlinegap0pt
\usepackage{bm}
\usepackage{bbm}
\usepackage{dsfont}
\usepackage{xspace}
\usepackage{cancel}
\usepackage{mathtools}
\usepackage{amsthm}
\makeatletter
\@ifundefined{thm}{\newtheorem{thm}{Thm}}{}
\@ifundefined{Def}{\newtheorem{Def}[thm]{Definition}}{}
\@ifundefined{cor}{\newtheorem{cor}[thm]{Corollary}}{}
\@ifundefined{lem}{\newtheorem{lem}[thm]{Lemma}}{}
\makeatother
\newcommand{\R}{\mathbb{R}}
\newcommand{\Z}{\mathbb{Z}}
\newcommand{\norm}[1]{\left\lVert#1\right\rVert}
\title[Calibrated submanifolds, adiabatic limit, and gradient graph]{Calibrated submanifolds, adiabatic limit, and gradient graphs}
\author{Yang Li}
\date{Source: arXiv:2608.19700v1 (CC BY 4.0)}
\begin{document}
\maketitle

\noindent\emph{Source and licence.} Calibrated submanifolds, adiabatic limit, and gradient graphs. Version arXiv:2608.19700v1 (CC BY 4.0), distributed under the Creative Commons Attribution 4.0 International licence, as recorded in the arXiv licence metadata for this exact version and shown on the abstract page https://arxiv.org/abs/2608.19700v1. Full rights notice: licenses/arxiv-2608.19700-CC-BY-4.0.txt.\par\smallskip

\noindent\textit{Editorial scope.} Counted unit: the Lem whose source label is \texttt{lem:integrality} in the article (source0.tex lines 1089--1091, source label \texttt{lem:integrality}), delivered together with the article's own complete original proof of it (source0.tex lines 1173--1281, one \texttt{proof} environment of 109 lines). No mathematics is written, repaired or re-derived: statement and proof are the article's own text line for line. The proof is detached from the statement in the source: the article states the result at lines 1089--1091 and prints its proof at lines 1173--1281, and the proof environment's own header names the statement, so the association is the article's own. Required context retained uncounted: Def:uniformvolumebound (lines 385--396); Def:gradientcycle (lines 503--514); cor:volumebound (lines 845--863). Every definition, display and label of those lines is retained unchanged. Omitted from this package and not redistributed: 1--384, 397--502, 515--844, 864--1088, 1092--1172, 1282--2119. Nothing inside a retained excerpt is omitted or altered: every retained body is delivered exactly as the article's lines are written, so no omission receipt of any kind accompanies this package. Citations: the retained excerpts carry no citation command at all, so no bibliography entry of the article is transcribed and no reference is redistributed. Reference adaptation: none was needed and none was made; every reference inside the delivered unit resolves inside it, so no unresolved reference or citation is produced. Printed slips kept literal: no printed word, symbol, hypothesis or number of the counted statement or of its proof was altered. Layout and class: the article's own class is \texttt{article} with packages including babel, inputenc, stmaryrd, amsmath, graphicx, amssymb, amsthm, tikz-cd, mathrsfs, todonotes, enumitem, yfonts, dsfont, MnSymbol; the wrapper is the lane's pinned \texttt{amsart} editorial wrapper at 10pt on A4, which restates the theorem-like environments the retained text needs and adds the article's own macro definitions that the retained text needs (\texttt{\textbackslash R}, \texttt{\textbackslash Z}, \texttt{\textbackslash norm}), with extra wrapper packages (bm, bbm, dsfont, xspace, cancel, mathtools). The delivered numbering is the unit's own. Assets: no retained span contains a figure environment, a graphics-inclusion command, a PDF-page inclusion, a table or a TikZ picture; no image file is copied into this package, the package declares no asset dependency and no image file is redistributed with it. Licence: version arXiv:2608.19700v1 of this article is distributed under the Creative Commons Attribution 4.0 International licence, as recorded in the arXiv licence metadata for this exact version and shown on the abstract page https://arxiv.org/abs/2608.19700v1. A full rights notice is delivered at licenses/arxiv-2608.19700-CC-BY-4.0.txt. Characters: every retained excerpt is pure ASCII, so the delivered engine, XeLaTeX with its default Latin Modern text font, reports no missing character.
	\begin{Def}\label{Def:uniformvolumebound}
	We say the special Lagrangians $L_i$ inside $(X, g_\epsilon, \omega_\epsilon, \Omega_\epsilon)$ (resp. associative $L_i$ inside $(X, g_\epsilon, \phi_\epsilon)$) satisfy the \emph{uniform linear	volume bound}, if $\epsilon \to 0 $ as $i\to +\infty$, and 
	for each given compact subset $K\subset B$, there is 
	some uniform constant $r_0, C$ depending on $K$ but independent of $i,\epsilon$, such that $B_g(p, r_0)$ is properly contained in $B$, and 
	\begin{equation}
\sup_{   p\in K }	\epsilon^{1-n} \norm{L_i} (  \pi^{-1} B_g(p, r_0)   ) \leq C.
	\end{equation}
	resp. 
	\[
	\sup_{p\in K }	\epsilon^{-2} \norm{L_i}( \pi^{-1} B_g(p, r_0)   ) \leq C.
	\]
	\end{Def}

\begin{Def}\label{Def:gradientcycle}
Let $T$ be a $\Lambda$-valued 1-rectifiable current. %We define the \emph{density function} $\Theta:\Gamma\to \R_{\geq 0}$ by
%\begin{equation}
%\Theta(x)=  |\delta(x) |_x.
%\end{equation}
%(As a caveat, the density may not be integer valued!) 
The \emph{density measure} for $T$ is the Radon measure 
\[
\mu_T (f)= \int_\Gamma f(x) |\delta(x)|_x d\mathcal{H}^1 (x),\quad \forall f\in C^0_{cpt}(B).
\]

\end{Def}

		\begin{cor}\label{cor:volumebound}
		
We assume the special Lagrangian (resp. associative) cycles $L$ satisfy the uniform linear mass bound in Def. \ref{Def:uniformvolumebound}. Then there is a uniform in $\epsilon$ mass upper bound for $\epsilon\leq r<r_0$,
		\[
		\norm{L}( \pi^{-1}B_g(p, r) ) \leq \begin{cases}
			C\epsilon^{n-1}r ,\quad &\text{special Lagrangian case},
			\\
		C	\epsilon^2 r,\quad  & \text{associative case}.
		\end{cases}
		\]
		Furthermore, if $\pi^{-1}(p)$ contains some point in the support of $L$, then there is a uniform mass lower bound for $\epsilon\leq r<r_0$,
		\[
		\norm{L}( \pi^{-1}B_g(p, r) ) \geq \begin{cases}
			C^{-1}\epsilon^{n-1} r ,\quad &\text{special Lagrangian case},
			\\
			C^{-1}	\epsilon^2r ,\quad  & \text{associative case}.
		\end{cases}
		\]
		\end{cor}

		\begin{lem}\label{lem:integrality}
		The class $\delta(p)$ lies in the integral lattice $\Lambda= H_{n-1}(T^n,\Z)$ (resp. $H_2(K3,\Z))$.
		\end{lem}

		\begin{proof}
			We focus on the associative case, and continue with the notations in the proof of Lemma \ref{lem:integrality}. Recall $0< \gamma\ll 1$ is any given small number.  We suppose $0<r<\gamma$, and $r$ is small enough as required in the arguments of Lemma \ref{lem:integrality}. The constants $C$ below are independent of $\gamma, \epsilon, i$. 
			
			
			
		
	%	Moreover, for $i$ large enough depending on $r$, we have $\epsilon <r$, hence by the mass upper bound in Cor. \ref{cor:volumebound}
		%	\[
		%	\text{Mass}(   L_i \lfloor  C_{p,r}     ) \leq  \norm{L_i}( \pi^{-1}(B_g(p, 10r))   )\leq C\epsilon^2 r. 
		%	\]
%Thus by the coarea formula, we can \emph{select a good slice} $r/2< t<r$, such that 
%			\begin{equation}\label{eqn:goodslice}
%				\begin{cases}
%				\text{Mass}   (  {L_{i,t}} ) + 	\text{Mass}   (  {L_{i,-t}} ) \leq C\epsilon^2,
	%			\\
%				\int_{ \text{spt}(L_{i,t}   )  }( 1- |\nabla d_p| ) d\norm{ L_{i,t}}  +    	\int_{ \text{spt}(L_{i,-t}   )  }( 1- |\nabla d_p| ) d\norm{ L_{i,-t}}    \leq C\epsilon^2 \gamma r
%				\end{cases}
%			\end{equation}
			
			
			
			
			
			
		
			
	By the coarea formula, and the mass upper bound in Cor. \ref{cor:volumebound}, when $i$ is large enough so that $\epsilon\ll r$,
			\[
			\int_{r/2}^r  \text{Mass}   (  {L_{i,t}} )  dt \leq C \text{Mass}(   L_i \lfloor  C_{p,r}     ) \leq C \norm{L_i}( \pi^{-1}(B_g(p, 9r))   )\leq C\epsilon^2 r. 
			\]
			Similarly
		$
				\int_{r/2}^r  \text{Mass}   (  {L_{i,-t}} )  dt \leq C\epsilon^2 r. 
$
			Thus we can select a  good slice $r/2<t<r$, such that 
			\begin{equation}
			\text{Mass}   (  {L_{i,t}} ) + 	\text{Mass}   (  {L_{i,-t}} ) \leq C\epsilon^2.
			\end{equation}
			
			
			
		
			
			
			

			
			
			Recall that $L_{i,t}\subset   \{   \text{dist}_g(x, p+ \R v) <2 \gamma r   \}   $ as in the proof of Lemma \ref{lem:integrality}. Since $\gamma \ll 1$ is small, $\partial ( L\lfloor C_{p,t} )= L_{i,t}- L_{i,-t}$ as there is no boundary contribution for $|t|<r$. We take local coordinates $x_1, x_2, x_3$ so that $\frac{\partial}{\partial x_i}$ is $g$-orthonormal at $p$, and $x_1$ agrees with $\pi_{p,v}$, so in particular $v= \xi(p)= \frac{\partial}{\partial x_1}$. 
			Then by Stokes formula,
			\begin{equation}
				\begin{split}
					&	\int_{L_i  \lfloor C_{p,t} }(\sum_i \epsilon^2 \omega_i \wedge dx_i - dx_1dx_2dx_3)
					\\
					= & \int_{L_{i,t}} \sum_{cyc}  x_1(\epsilon^2\omega_1- \frac{1}{3}dx_2\wedge dx_3) -   \int_{L_{i,-t}} \sum_{cyc}  x_1(\epsilon^2\omega_1- \frac{1}{3}dx_2\wedge dx_3) .
				\end{split}
			\end{equation}
			Since $\phi_\epsilon$ deviates from the product structure by $O(r)$ relative error, and $r<\gamma$,
			\[
				\int_{L_i   \lfloor C_{p,t}   } \sum_i \epsilon^2\omega_i \wedge dx_i - dx_1dx_2dx_3= \int_{L_i   \lfloor C_{p,t}   } \phi_\epsilon +O(r\norm{L_i}(  C_{i,t}   ) )= \norm{L_i}(  C_{p,t}   ) (1+O(\gamma)).
			\]
			Since $L_{i,t}\subset   \{   x_1=t,  \text{dist}_g(x, p+ \R v) <2 \gamma r   \}   $, we know $|x_2|+ |x_3| \leq C\gamma r \leq C\gamma t$ on the support of $L_{i,t}$, so
			\[
			\begin{split}
					\int_{L_{i,t}} \sum_{cyc}  x_1(\epsilon^2 \omega_1- \frac{1}{3} dx_2\wedge dx_3) =&  t	\int_{L_{i,t}}  (\epsilon^2 \omega_1- \frac{1}{3} dx_2\wedge dx_3)  + O(\gamma  t \text{Mass}   (  {L_{i,t}} )  )
					\\
					=&  t	\int_{L_{i,t}}  (\epsilon^2 \omega_1- \frac{1}{3} dx_2\wedge dx_3)  + O(\gamma  t \epsilon^2 )
					\\
					= & t \epsilon^2 \langle \delta(p) , [\omega_1] \rangle + O(\gamma  t \epsilon^2 )
			\end{split}
			\]
			where the second line uses the mass bound for the good slice, and the third line evaluates the cohomological integral on the closed current $L_{i,t}$, using $[L_{i,t}]= \delta_r=\delta(p)\in H_2(K3,\Z)$ for small enough $r$. Likewise with $L_{i,-t}$.
			
			
			
			Combining these ingredients in the above Stokes formula,
			\begin{equation}\label{eqn:massincylinder}
			\norm{L_i}(  C_{p,t}   )= 2\epsilon^2 t (1+ O(\gamma ) ) \langle \delta(p) , [\omega_1] \rangle + O(\gamma  t \epsilon^2) .
			\end{equation}
		Since $|d\pi_{p,v}|\leq 1+Cr\leq 1+C\gamma$, the slicing inequality gives
				\[
				\int_{-t}^{t}\operatorname{Mass}(L_{i,s})\,ds
				\leq (1+C\gamma)\,\norm{L_i}(C_{p,t}).
				\]
				There is no factor $2t$ in this coarea estimate. Dividing by the length $2t$ of the slicing interval only when taking the average, and then using \eqref{eqn:massincylinder}, we obtain a slice with $|s|<t<r$ for which
				\[
				\begin{split}
					\operatorname{Mass}_{g_\epsilon}(L_{i,s})
					&\leq \frac{1+C\gamma}{2t}\,\norm{L_i}(C_{p,t})\\
					&\leq (1+C\gamma)\epsilon^2\langle\delta(p),[\omega_1]\rangle
					+C\gamma\epsilon^2.
				\end{split}
				\]
			We can use a $(1+C\gamma)$-Lipschitz map to project the closed integral current $L_{i,s}$ into the K3 fibre $\pi^{-1}(p)$ with metric $\epsilon^2 g_{K3}$. This gives a representative of $\delta(p)\in H_2(\pi^{-1}(p),\Z )$, with $g_{K3}$-mass bounded by $(1+ C\gamma)\langle \delta(p) , [\omega_1] \rangle + C\gamma $, hence
			\begin{equation*}
	|\delta(p) |_p \leq (1+ C\gamma)\langle \delta(p) , [\omega_1] \rangle + C\gamma .
			\end{equation*}
			This holds for any given small $\gamma>0$, hence 
$
	|\delta(p)|_p \leq  \langle \delta(p) , [\omega_1] \rangle .
$
Thus the mass minimizer of $\delta(p) \in H_2(\pi^{-1}(p),\Z )$ must be calibrated by the K\"ahler form $\omega_1$, so it is an $I_1$-holomorphic curve (possibly with multiplicity). In particular, for $v= \xi(p)= \frac{\partial}{\partial x_1}$, 
\begin{equation}\label{eqn:densitycalibration}
|\delta(p)|_p = \langle \delta(p) , [\omega_1] \rangle = \langle \delta(p), \nabla_v H\rangle ,
\end{equation}
and $\int_\delta [\omega_2]= \int_\delta [\omega_3]=0$, so that $\langle \delta(p), \nabla H\rangle $ points in the direction of $v$. This verifies the gradient flow condition in Def. \ref{Def:gradientcycle}. We recall that $\partial T=0$, hence $T$ is a $\Lambda$-weighted gradient cycle.

The special Lagrangian case follows the same strategy up to cosmetic changes.
		\end{proof}

\end{document}
